\section{Detection and Tracking Application Evaluation}
\label{sec:tracking_app}
The proposed detection and tracking subsystem is responsible for receiving the processed 
radar data
across all radar channels from the aggregator subsystem, and producing tracks of
targets in the radar's field of view. The subsystem performs the following operations
sequentially (see Figure \ref{fig:dat-dataflow}): CFAR, clustering, monopulse 
angle sensing, sidelobe suppression (blanking), and multi-target tracking.
Below are the implementation and experiment details for each of the operations.

\begin{figure}[H]
    \centering
    \incfig[1]{detection_and_tracking_dataflow} 
    \caption[Dataflow diagram for the detection and tracking subsystem.]{Dataflow for the detection 
    and tracking subsystem. The subsystem receives the processed radar data across all radar 
    channels from the aggregator subsystem, and produces tracks of targets in the radar's 
    field of view. The subsystem performs the following operations sequentially: CFAR, clustering, 
    monopulse angle sensing, sidelobe suppression (blanking), and multi-target tracking.}
    \label{fig:dat-dataflow}
\end{figure}

\subsection{CFAR}
For finding the CFAR coefficient $\alpha$ (defined in \ref{eqn:threshold}) that gives a desired false 
alarm rate $P_{fa}^{\star}$. The relationship between these variables (Equation 
\ref{eqn:cfar_false_alarm}), cannot be rearranged into a closed-form expression. 

\begin{equation}
    P_{fa}(\alpha;N) = \left( \frac{N}{N+\alpha} \right)^N
    \sum_{k=0}^{3} \binom{N+k-1}{k}
    \left( \frac{\alpha}{N+\alpha} \right)^k
\end{equation}

Since $P_{fa}(\alpha;N)$ is strictly decreasing in $\alpha$, bisection can be used 
to find the value of $\alpha$ that gives the desired false alarm rate \(P_{fa}^{\star}\).

The method begins with $\alpha_{lo}=0$ and increases an upper bound $\alpha_{hi}$ until 
$P_{fa}(\alpha_{hi};N) \leq P_{fa}^{\star}$. It then repeatedly bisects the interval and evaluates 
$P_{fa}(\alpha;N)$ at the midpoint, updating the lower or upper bound depending on whether the resulting 
false alarm rate is above or below the target. This continues until the interval is sufficiently small. 
In this way, $\alpha$ is obtained as the numerical solution that gives the required false alarm rate.

Additionally, three different implementations of CFAR are compared for their latency: a naive
implementation on the CPU, an implementation that uses a prefix sum table, and a naive solution on the
GPU. Additionally, different noise assumptions are tested. The naive implementation involves calculating 
the noise statistic with a repeated sum per sample of interest (see Figure \ref{fig:naive-cfar}). 
The prefix sum table method reduces the 
number of sums required to four per sample of interest regardless of the number of training cells, the
start and end of the left and right window (see 
Figure \ref{fig:prefix-sum-table}). The GPU implementation creates a thread per sample of interest (see 
Figure \ref{fig:gpu-cfar}). A sweep of window size, guard size, frame buffer size and false alarm rate values is 
conducted for each CFAR implementation on real and simulated data. The performance (throughput, latency, 
false alarm rate, sensitivity, and track quality) is compared across configurations in a cost benefit 
analysis. 
This addresses the literature gap in understanding how detection sensitivity and false alarm behaviour 
trade off against latency and throughput under SWaP constraints.

\begin{figure}[!ht]
    \incfig[0.75]{cfar}
    \caption[Naive CPU CFAR implementation.]{Naive Implementation of CFAR for the CPU. Each sample is compared to a threshold
    determined by the power of the samples around it. The threshold is calculated by taking the
    average power of the surrounding samples and multiplying it by a factor \(\alpha\) determined
    by the desired false alarm rate.}
    \label{fig:naive-cfar}
\end{figure}

\begin{figure}[htbp]
    \incfig[0.75]{prefix_sum_table}
    \label{fig:prefix-sum-table}
    \caption[CFAR using a prefix sum table for constant-time noise statistics.]{CFAR implementation using a prefix sum table (cumulative signal power). The prefix sum
    table allows the noise statistic to be calculated with a fixed number of sums per sample of
    interest, regardless of the number of training cells. In this specific case, after the cumulative
    signal power is calculated in processing, it only takes four sums of to find the sum of eight
    training cells.}
\end{figure}

\begin{figure}[htbp]
    \incfig[0.75]{gpu_cfar}
    \label{fig:gpu-cfar}
    \caption[Parallel GPU CFAR implementation.]{CFAR implementation on the GPU. Threads (each depicted as unique colours) are each
    responsible for a sample, allowing for parallel processing of the CFAR algorithm across the entire
    range profile. Each thread calculates its own noise statistic, threshold, and mask.}
\end{figure}

The literature gap regarding algorithmic selection is addressed through an exploration of different CFAR
algorithms, including OS-CFAR, CA-CFAR, SO-CFAR, and GO-CFAR. The effect of these algorithms on detection
sensitivity and track quality is compared. For intuition, Figure \ref{fig:cfar-algorithms} illustrates
the difference in performance of the algorithms on a synthetic signal with 5 peaks. This work 
compares the performance of GO-CFAR, SO-CFAR, and CA-CFAR on simulated data and discusses
the downstream effects of this selection, addressing the literature gap regarding the impact of 
detection algorithms in the noise regimes of miniature radar. Additionally, by analysing the 
performance of both CPU and GPU implementations of CFAR, the literature gap regarding the 
algorithmic implementation under SWaP constraints is addressed. 

\begin{figure}
    \centering
    \incfig[0.81]{cfar_stacked_matched_colours}
    \caption{ Performance of four CFAR variants. Each 
    algorithm registers a detection when its threshold is 
    below the signal. The test signal has 5 peaks that are 
    each surrounded by different noisy distributions. The bottom 
    plot illustrates the detections made by each algorithm. 
    OS-CFAR performs well in this example.}
    \label{fig:cfar-algorithms}
\end{figure}

\subsection{Clustering}
After CFAR occurs, most of the samples (range bins) will be zero. The efficiency loss from 
processing so many zero samples motivates a stage that transforms the detection sparse 
representation into a dense one. Additionally, the widening of the main lobe of each of 
the detections due to the windowing stage creates multiple range bins with detections 
from the same target, which motivates a stage that groups these samples into a single 
detection.
The algorithm that addresses both of these problems, clustering, iterates through the 
range profile, and creates a cluster for each contiguous set of 
detection samples. Then, the complex signals from each of the samples are summed together. 
SNR is expected to improve, as noise will be incoherent within clusters, but the signal
that corresponds to a target will have the same phase across samples, and thus will be 
coherent. The latency of both CPU and GPU implementations of the clustering algorithm 
is measured. Different clustering parameters are tested for their impact on tracking performance, such 
as the allowed gap size between detections.

\subsection{Monopulse Angle Sensing}
Angle sensing converts multi-channel detections into azimuth, elevation and range estimates that 
drive association and tracking. We assess candidate angle discriminators by running controlled 
sweeps over SNR, channel 
imbalance, and multi-target proximity, and measuring angular bias, mean squared error, and confidence. 
The latency and throughput of CPU and GPU implementations are measured. The trade-off involves software implementation details such as warp divergence and instruction-level parallelism. 

The angle sensing algorithm features a few early returns that result in warp divergence on the GPU. 
For example, if the bearing and elevation ratios are outside of the domain over which the polynomial
fit is valid. Instruction level parallelism is a measure of how many instructions can be executed 
in parallel on a single thread. By changing the order of operations, we can potentially increase instruction-level parallelism and reduce the latency of the GPU implementation.

\subsection{Sidelobe Suppression (Blanking)}
When a target exists outside of the field of view of a given radar receiver, the receiver may still
register a detection due to the sidelobes of the antenna pattern. When a target is in the radar's 
field of view, it will land in the main lobe of some receivers and the sidelobes of others (see Figure 
\ref{fig:sidelobes}), resulting
in a strong detection in the correct direction and weaker detections in the wrong directions, all at
the same range. The blanking stage eliminates detections that are likely to be sidelobes by eliminating
detections that are close in range to a stronger detection at the same time step. 

\begin{figure}[htbp]
    \centering
    \begin{minipage}{0.5\textwidth}
        \centering
        \incfig[1]{sidelobes}
        \caption[Sidelobes of the antenna pattern.]{Receivers R1 and R3, which are pointing away from the target,
        still receive significant power from the first sidelobes lining up with the target. }
        \label{fig:sidelobes}
    \end{minipage}
    \begin{minipage}{0.45\textwidth}
        \centering
        \incfig[1]{blanking}
        \caption[Blanking for sidelobe suppression.]{Detections from the sidelobes of R1 and R3 are weaker 
        than R2's detection from the main lobe, and are close in range to R2's detection. R1 and R3's
        detections that are close in range to a stronger detection at the same time step are eliminated.}
        \label{fig:blanking}
    \end{minipage}
\end{figure}

\begin{figure}[htbp]
\end{figure}

The blanking stage will be implemented on the CPU. After clustering and angle sensing, 
the detections come in an ordered list per receiver quadruplet. A CPU implementation would involve,
for each time step, iterating through each of the lists simultaneously and sorting them into one ordered list in O(N) 
time. Then iterating through the ordered list and eliminating detections that are close in range to
a stronger detection. 

% GPU DETAILS BELOW, POTENTIALLY USEFUL
% A GPU implementation would involve creating a thread per detection per quadruplet, 
% performing 3 binary searches to figure out how many detections are nearer than it in range to the radar,
% and mapping itself to the corresponding place in the fully sorted list. In a second pass then, each 
% thread would check the detections near it in the fully sorted list to see if there are any stronger
% detections in range. Since this problem is a standard k-way merge sort, the CUDA Thrust library 
% implementation will be used. This stage can also double as a receiver fusion stage, as detections
% from multiple receivers across different chirps that are close in not just range but also angular 
% position can be merged 
% together. 

\subsection{Multi-Target Tracking}

\begin{figure}[htbp]
    \centering
    \incfig[1]{tracking_dataflow} 
    \caption[Dataflow diagram for the tracking subsystem.]{Dataflow for the tracking stage of the 
    detection and tracking subsystem. Detections are fed into the
    associator of the multi-target tracker, and are then fed into the hypothesiser. 
    The hypothesiser produces a hypothesis per track/detection pair. The associator then 
    finds the optimal assignment
    of detections to tracks. The updater then updates the tracks, and the deleter deletes 
    tracks with poor quality. For the detections that are not associated with any tracks,
    a tentative track is initiated or updated in the initiator, a miniature tracker that
    handles the track initiation and confirmation process. 
    }
    \label{fig:tracking-dataflow}
\end{figure}

The multi-target tracking stage will be implemented entirely on the CPU due to its sequential nature. 
Amplitude based association will be compared to non-amplitude based association for their impact on
track quality. The multi-target tracker has over 20 parameters that can be tuned, and one 
configuration of these parameters will be validated and discussed. Experiments will
be conducted in a variation of SNR, and number of targets, in simulated data. Real 
data will be used to validate the results on static targets. The literature gap regarding the impact
of algorithmic choices in the tracking stage on track quality is addressed through this experiment. 

% EXPLORE MORE
% - Non amplitude based association
% - Amplitude based association
% - ringing 
